\PassOptionsToPackage{dvipsnames}{xcolor} % colori PDF/A
\PassOptionsToPackage{hyperfootnotes=false}{hyperref} % note a piè di pagina non cliccabili nel testo (v. anche punto 7)

\usepackage{colorprofiles}
% PDF/A
% validate in https://www.pdf-online.com/osa/validate.aspx
\usepackage[a-1a,mathxmp]{pdfx}[2018/12/22]
\usepackage[T1]{fontenc}
\usepackage[utf8]{inputenc}
\usepackage[italian]{babel}
\usepackage{bookmark}
\usepackage{caption}
\usepackage{comment}
\usepackage{chngpage, calc} % centra il frontespizio
\usepackage{emptypage} % pagine vuote senza testatina e piede di pagina
\usepackage{epigraph} % per epigrafi
\usepackage{graphicx} % immagini
\usepackage[pdfa]{hyperref} % collegamenti ipertestuali
\usepackage{mparhack,relsize}  % finezze tipografiche
\usepackage{nameref} % visualizza nome dei riferimenti
\usepackage[font=small]{quoting} % citazioni
\usepackage{subfig} % sottofigure, sottotabelle
\usepackage[italian]{varioref} % riferimenti completi della pagina
\usepackage{longtable} % tabelle su più pagine
% NOTA SU "automake": il pacchetto "glossaries" offrirebbe un'opzione
% "automake" (o "automake=lite") per rigenerare da sé, a fine compilazione,
% gli elenchi di Acronimi e abbreviazioni/Glossario (i file .gls/.acr,
% ottenuti processando i file grezzi .glo/.acn). Questa opzione NON è però
% usata qui, perché richiede che il motore TeX possa eseguire un comando di
% sistema (\write18) per "makeglossaries"/"makeglossaries-lite": MiKTeX, per
% impostazione predefinita, esegue il flag -shell-escape in modalità
% "ristretta" (restricted), che consente solo un elenco prestabilito di
% programmi (bibtex, biber, makeindex, ecc.) e blocca qualunque altro
% comando, incluso "makeglossaries"/"makeglossaries-lite" (nel log MiKTeX
% questo compare come "runsystem(makeglossaries-lite ...)...disabled
% (restricted)."). Il risultato è che gli elenchi restano vuoti in stampa e
% i relativi collegamenti ipertestuali (\gls, ecc.) puntano a un'ancora mai
% creata, risultando "a vuoto" una volta cliccati.
% La rigenerazione degli elenchi è quindi affidata non a TeX ma a latexmk
% stesso (v. il file "latexmkrc" nella stessa cartella di questo file), che
% lancia "makeglossaries" come processo esterno al di fuori del motore TeX,
% aggirando così del tutto la restrizione dello shell escape di MiKTeX.
\usepackage[toc, acronym]{glossaries}
\usepackage[backend=biber, style=ieee, hyperref=false, backref=true, defernumbers=true]{biblatex} % stile IEEE per le fonti non giuridiche; citazioni non cliccabili nel testo (v. thesis_config.tex). Le fonti giuridiche restano escluse da questo stile e sono gestite da \legfootcite (v. thesis_config.tex). defernumbers=true: la bibliografia è divisa in più sezioni (Bibliografia: Testi/Articoli/Normativa; Sitografia: Siti), quindi la numerazione IEEE va calcolata solo alla fine per restare coerente tra le sezioni
\usepackage{lmodern}
\usepackage[top=3cm, bottom=3cm, inner=3cm, outer=2.5cm]{geometry} % margini specchiati per la rilegatura (richiede twoside nella \documentclass)
\usepackage{parskip} % paragrafi separati da spazio verticale, senza rientro di prima riga
\usepackage{fancyhdr}
\usepackage{lipsum}
\usepackage{setspace}
\usepackage{titlesec}
\usepackage[cachedir=minted-caches]{minted} % https://it.overleaf.com/learn/latex/Code_Highlighting_with_minted
\usepackage{xcolor}
\usepackage{newunicodechar} % permette di "disinnescare" caratteri Unicode invisibili incollati per errore (vedi thesis_config.tex)
\usepackage{csquotes} % gestisce automaticamente i caratteri (")
\usepackage{etoolbox}
\usepackage[bottom]{footmisc}
\usepackage{zref-totpages}